\documentclass[10pt,a4paper]{amsart}
\usepackage[margin=19mm]{geometry}
\usepackage{amsmath,amssymb,mathtools}
\usepackage[scr=rsfs,cal=euler]{mathalfa}
\usepackage{xfrac}
\usepackage[unicode,hidelinks,bookmarks=false]{hyperref}
\newcommand{\ie}{i.e.\ }
\newcommand{\cf}{cf.\ }
\newcommand{\resp}{respectively\ }
\newcommand{\Subsection}{\subsection}
\providecommand{\nobreakdash}{\nobreak\hbox{-}}
\setlength{\emergencystretch}{2em}
\multlinegap0pt
\usepackage{bm}
\usepackage{amssymb}
\usepackage{xcolor}
\usepackage{tikz}
\usepackage[normalem]{ulem}
\newtheorem{theorem}{Theorem}
\title{Piecewise Linear Functions and Neural Network Expressivity via Discriminantal Arrangements}
\author{Pragnya Das}
\date{Source: arXiv:2604.02480v1 (CC BY 4.0)}
\begin{document}
\maketitle

\noindent\emph{Source and licence.} Piecewise Linear Functions and Neural Network Expressivity via Discriminantal Arrangements. Version arXiv:2604.02480v1 (CC BY 4.0), distributed by arXiv under the Creative Commons Attribution 4.0 International licence, http://creativecommons.org/licenses/by/4.0/, as recorded in the arXiv licence metadata for this exact version and shown on the abstract page https://arxiv.org/abs/2604.02480v1. Full licence notice: \texttt{licenses/arxiv-2604.02480-CC-BY-4.0.txt}.\par\smallskip

\noindent\textit{Editorial scope.} Counted unit: the article's theorem thm:indicator (source0.tex lines 272--286). The article's complete original proof of it is delivered with it (source0.tex lines 287--326, one \texttt{proof} environment of 40 lines). No mathematics is written, repaired or re-derived: the statement and the proof are the article's own lines, in the article's own notation, and no retained line is substituted or re-flowed.  No further source-local context is needed: every cross-reference of every retained line resolves inside the two delivered spans.  Nothing was omitted from any retained span, so the package carries no omission record.  No retained line contains a citation, so the wrapper carries no bibliography and no bibliography entry.  No retained line contains a figure environment, an image-inclusion command or drawing code, so no image is delivered and the package declares no asset dependency.  Editorial formatting only, in addition: an AMS \texttt{amsart} preamble that restates the article's own declarations and macros used by the retained lines, namely the environments \texttt{theorem} on separate counters, together with the standard packages those lines need.  The retained environments print in this document as Theorem 1 in order of appearance, whereas the article prints the same items with the numbering of its own full document.  The complete native source of the article as distributed in the arXiv e-print for this exact version is bundled unchanged as \texttt{sources/197739/source0.tex}.
\begin{theorem}\label{thm:indicator}
Let $A$ be a hyperplane arrangement in $\mathbb{R}^n$ with simplicial fan $\Sigma(A)$. Assume that for every full-dimensional cone $\sigma \in \Sigma(A)$, the set
\[
V_\sigma := \{1_S \in \sigma : S \subseteq [n]\}
\]
affinely spans $\mathbb{R}^n$. Then every function $F \in \mathcal{F}_A$ admits a realization by a continuous piecewise linear function $f$ compatible with $A$, i.e.,
\[
f(1_S) = F(S) \quad \text{for all } S \subseteq [n].
\]
Equivalently, under this condition, the map
\[
\mathcal{V}_A \to \mathcal{F}_A, \quad f \mapsto F
\]
is surjective.
\end{theorem}

\begin{proof}
Let $F \in \mathcal{F}_{\mathcal{A}}$ be given.\\
\noindent
\textbf{Step 1: Local affine interpolation.}
Fix a full-dimensional cone $\sigma \in \Sigma(\mathcal{A})$. By assumption, the set $V_\sigma$ affinely spans $\mathbb{R}^n$. Hence there exists a unique affine function $f_\sigma : \mathbb{R}^n \to \mathbb{R}$ such that
\[
f_\sigma(\mathbf{1}_S) = F(S) \quad \text{for all } \mathbf{1}_S \in V_\sigma.
\]
\noindent
\textbf{Step 2: Consistency on overlaps.}
Let $\sigma, \tau \in \Sigma(\mathcal{A})$ be two cones sharing a common face $\sigma \cap \tau$. \\
\noindent
For any indicator vector $\mathbf{1}_S \in \sigma \cap \tau$, we have
\[
f_\sigma(\mathbf{1}_S) = F(S) = f_\tau(\mathbf{1}_S).
\]
\noindent
Since both $f_\sigma$ and $f_\tau$ are affine functions and agree on a set of points whose affine span equals $\operatorname{aff}(\sigma \cap \tau)$, it follows that
\[
f_\sigma = f_\tau \quad \text{on } \sigma \cap \tau.
\]
\noindent
\textbf{Step 3: Global definition.}
Define a function $f : \mathbb{R}^n \to \mathbb{R}$ by
\[
f(x) := f_\sigma(x) \quad \text{if } x \in \sigma.
\]
\noindent
This is well-defined by Step 2.\\
\noindent
\textbf{Step 4: Properties.}
By construction:
\begin{itemize}
\item $f$ is affine on each cone $\sigma \in \Sigma(\mathcal{A})$,
\item $f$ is continuous across cone boundaries,
\item $f(\mathbf{1}_S) = F(S)$ for all $S \subseteq [n]$.
\end{itemize}

Thus $f \in \mathcal{V}_{\mathcal{A}}$, and the map $f \mapsto F$ is surjective.
\end{proof}

\end{document}
